\section{The two orientation counts}
\label{sec:joint}
Write \(X_n^+,X_n^-\) for the positive and negative counts. For \(u,v\ge0\) let
\[
 Z_n(u,v)=\sum_{P\in\M_n}u^{X_n^+(P)}v^{X_n^-(P)},\qquad
 f_{u,v}(z)=z\left(\frac1{1-(u-1)z}+
                        \frac1{1-(v-1)z}-1\right).
\]
Here \(\M_n\) is the set of connected diagrams, and \(0^0=1\).
Let \(R=1/\mu\) and \(\M(z)=\sum_{n\ge1}M_nz^n\).

\begin{theorem}\label{thm:joint}
As a formal power series,
\[
             \sum_{n\ge1} Z_n(u,v)z^n=\M(f_{u,v}(z)).
\]
For \(u,v>0\), let \(\rho(u,v)\) be the first positive solution of
\(f_{u,v}(\rho)=R\), before any positive pole. Then
\(\lim Z_n(u,v)^{1/n}=1/\rho(u,v)\). Under the uniform measure,
\((X_n^+/n,X_n^-/n)\) satisfies a large-deviation principle at speed \(n\)
with good rate function
\[
 I(a,b)=\sup_{s,t\in\mathbb R}
       \bigl\{as+bt+\log\bigl(\mu\,\rho(e^s,e^t)\bigr)\bigr\}.
\]
\end{theorem}
\begin{proof}
Give each marked positive site weight \(u-1\), each marked negative
site weight \(v-1\), and sum over subsets of actual occurrences.
On contracting the marked runs, every even vertex of a diagram of
order \(m\) has three possibilities: it remains unchanged, carries a
nonempty positive run, or carries a nonempty negative run. A run of
length \(j\) increases the order by \(j\). The \(m\) vertices can be
expanded independently by Lemma~\ref{lem:contract}, which gives the
formal identity with both activities.

In the upper-arch splice of two diagrams the only positive occurrence
which can be lost is the first site of the right diagram. A negative
occurrence can be gained only at the seam. Indeed every short deleted
upper arch starts at an even label, the long new arch has length at
least three, and the new short upper arch starts at an odd label.
The lower matching remains unchanged. Removing the two crossing upper
arches inverts the splice at its specified position. Thus, for
\(u,v>0\),
\[
 Z_{r+s}(u,v)\ \geq\ \min(1,u^{-1})\min(1,v)\,Z_r(u,v)Z_s(u,v).
\]
Fekete's lemma, applied after this constant rescaling, gives root
convergence. The counts have an exponential upper bound by the
Catalan-square estimate and the bounded number of sites.

For nonnegative activities and real \(z\) before the first pole,
\[
 f'_{u,v}(z)=\frac1{(1-(u-1)z)^2}+\frac1{(1-(v-1)z)^2}-1
             \ \geq\ 2/(1+z)^2-1.
\]
On \([0,\sqrt2-1)\) the right-hand side is positive. If a positive pole
occurs within this interval the function tends to infinity there.
Otherwise \(f_{u,v}\ge f_{0,0}=z(1-z)/(1+z)\) and the latter reaches
\(3-2\sqrt2>R\) at its maximum, since \(\mu>11\) \cite{AP}. This proves
existence and uniqueness of the first root on the increasing branch,
with strictly positive derivative. The composition is analytic along
the positive interval before this root. Its coefficients are
nonnegative, so Pringsheim's theorem rules out a smaller radius.
Analyticity across the root would, by the local inverse, continue
\(\M\) across its positive singularity at \(R\). Hence the radius is
exactly \(\rho\). This radius argument includes zero activities.

The limiting logarithmic moment generating function for the two
counts is \(-\log(\mu\rho(e^s,e^t))\). It is differentiable on all of
\(\mathbb R^2\) by the implicit function theorem. The rescaled counts
lie in a compact triangle. The G\"artner--Ellis theorem
\cite[Theorem 2.3.6]{DZ} now gives the asserted full large-deviation
principle for open and closed sets.
\end{proof}

\begin{corollary}\label{cor:double}
Let \(A_n=Z_n(0,0)\) count diagrams avoiding both signs. Then
\[
 \lim_{n\to\infty} A_n^{1/n}
       = \frac{\mu-1+\sqrt{\mu^2-6\mu+1}}{2}.
\]
\end{corollary}
\begin{proof}
In a connected left factor of order \(r\ge2\), its last pair of
vertices cannot have both an upper and a lower short arch: that would
be an isolated cycle. If the lower arch is absent, use the upper
splice, so no negative site appears at the seam. Otherwise use the
lower splice; the absence of the last upper arch prevents a positive
site at the seam. The chosen colour is recovered from the two arches
crossing the cut. Thus \(A_{r+s}\ge A_r A_s\) for \(r\ge2,s\ge1\).
Here \(A_1=1,A_2=0,A_3=2\), so every \(A_n\) for \(n\ge3\) is
positive. For fixed \(k\ge3\) write \(n=jk+\ell\) with
\(3\le\ell\le k+2\); this proves root convergence by the usual
supermultiplicative argument. Its limit is the reciprocal of the
first solution of \(z(1-z)/(1+z)=1/\mu\).
\end{proof}
